\chapter{Why RealQM Needs No Relativity}\label{ch:relativity}
\index{relativity}\index{Dirac equation}\index{special relativity}\index{Lévy-Leblond theorem}\index{fine structure}\index{QED}\index{charges as configurations}

\begin{quote}
\itshape Relativity is sufficient for much of what the atom does; the question is whether it is necessary.\\
\normalfont\hfill --- working note, 2026.
\end{quote}
\bigskip

In the standard story of twentieth-century physics, relativity is not optional. The electron is said to require
the Dirac equation; its magnetic factor $g=2$, its spin, and the fine structure of spectra are presented as
relativistic; light and matter are joined only through the relativistic quantum field theory of QED; and the
whole edifice is built on Lorentz invariance and Minkowski spacetime. Against this, RealQM is a frankly
\emph{non-relativistic} theory --- Coulomb interactions and a Schr\"odinger-type variational principle on
ordinary three-dimensional space --- and yet the preceding chapters have used it to compute atomic energies
across the periodic table, covalent and hydrogen bonds, condensed phases, nuclear binding, and orbital
magnetism. This chapter asks the obvious question the success raises: if so much of the physics of matter follows
without relativity, what exactly was relativity doing, and where is it genuinely needed?

\section{Where standard quantum mechanics invokes relativity}\label{sec:where-relativity}

It is worth being specific about the places the standard account reaches for relativity.
\begin{itemize}
\item \textbf{The electron equation.} Dirac's 1928 equation is relativistic by construction, and is taken as the
correct fundamental description of the electron, the Schr\"odinger equation being its low-velocity limit.
\item \textbf{Spin and $g=2$.} The two-valued spin and its gyromagnetic factor $g=2$ fall out of the Dirac
equation, and are routinely described as relativistic in origin.
\item \textbf{Fine structure.} The spin--orbit coupling $\xi(r)\,\mathbf L\!\cdot\!\mathbf S$ and the associated
splitting of spectral lines are $v^2/c^2$ corrections, genuinely relativistic in \emph{magnitude}.
\item \textbf{Antimatter and pair creation.} The positron and the creation and annihilation of particles are
consequences of the relativistic equation and have no non-relativistic counterpart.
\item \textbf{Light.} The electromagnetic field is Lorentz invariant, and its full quantum coupling to matter ---
QED --- is a relativistic field theory.
\end{itemize}
The first two of these, we will argue, are not what they seem; the last three mark the genuine boundary.

\section{RealQM builds the structure of matter without it}\label{sec:without-relativity}

The bulk of the physics of matter --- everything Parts~I--IV of this book compute --- is the electrostatics of
real charge, and needs no relativity. Atomic total energies, the covalent bond and its helium limit, hydride
atomization energies, dimer geometries, condensed phases, and even nuclear binding read as protons and electrons
bound by Coulomb alone, all follow from a non-relativistic variational principle. This is not an approximation
awaiting relativistic correction; for the structure, energetics, and chemistry of \emph{light} atoms --- and for
the gross total energies of all of them --- it is the physics itself, and relativity does no work. One
qualification must be entered, and is taken up in \S\ref{sec:heavy-atoms}: the relativistic \emph{kinetic energy}
of the tightly bound inner charge in \emph{heavy} elements, which their chemistry (gold, mercury, lead)
demonstrably requires.

\section{Spin and \texorpdfstring{$g=2$}{g=2} are geometric, not relativistic}\label{sec:g2-nonrel}

The single most-cited reason to believe the electron \emph{requires} relativity is $g=2$: since it drops out of
the Dirac equation, it is taken to be a relativistic quantity. It is not. As shown in
Chapter~\ref{ch:magnetism-spin}, Lévy-Leblond's theorem (1967) derives $g=2$ from a first-order,
\emph{non-relativistic}, Galilean spinor equation: relativity is \emph{sufficient} to produce $g=2$ but not
\emph{necessary}, because the factor of two does not come from Lorentz invariance in the first place. It comes
from the SU(2) double cover of the rotation group --- a spinor returns to itself only after $720^\circ$ --- a
geometric fact about how orientations of a real extended object compose, scale-free and equally true in a
Galilean world. Writing the kinetic term with $\boldsymbol\sigma\!\cdot\!\mathbf p$ rather than $p^2$ produces
the $g=2$ Zeeman coupling automatically, with no relativity used. So the phenomenon most responsible for the
belief that the electron is inescapably relativistic is, on inspection, a piece of non-relativistic geometry ---
exactly the residue RealQM aims to derive rather than import.

\section{Charges as configurations, not relativistic objects}\label{sec:charges-config}

Underlying the standard picture is the idea of the electron as a point object with an intrinsic rest mass, a thing
whose very kinematics is relativistic. RealQM offers a different reading: elementary charges are
\emph{configurations} of charge density in space, held together by the same Coulomb energetics as everything else,
their apparent mass and rigidity emergent properties of the configuration rather than primitive relativistic
attributes. On this reading there is no point electron whose motion demands Lorentz kinematics; there is an
extended charge distribution obeying continuum mechanics, and the ``particle'' is a stable pattern in it. The
need for relativity as the \emph{kinematics of point particles} dissolves along with the point particle.

\section{Light, time, and the reconstruction of relativity}\label{sec:light-time}

Two honest qualifications must be entered here, and they concern exactly the places relativity is real.

\emph{Light.} Maxwell's equations \emph{are} Lorentz invariant, and light is a relativistic object. But RealQM
does not need to become relativistic to handle it: as Chapter~\ref{ch:sm-coupling} showed, a real charge density
couples directly to a classical Maxwell field as a source, with no configuration-space wave function to promote to
a quantum field. The charge stays non-relativistic; the field is Maxwell's; the elaborate relativistic machinery
of QED, needed in the standard theory precisely because its matter lives on $3N$-space, is not required to join
real charge to real light.

\emph{Time and spacetime.} RealQM does not presuppose a Minkowski background. As Chapter~\ref{ch:time} argued,
the quiet atom is timeless and time is relational, entering only with motion; there is no spacetime manifold the
state is threaded through. And relativity itself, in the companion volume \emph{Many-Minds Relativity}, is
\emph{reconstructed} constructively --- each observer carrying an individual coordinate system, the Lorentz
transformations emerging from concrete signal-exchange protocols --- rather than assumed as a fixed arena. Where
the standard theory takes relativistic spacetime as a given backdrop, RealQM either does without it (bound matter)
or builds it from observation (Vol.~VII).

\section{Heavy atoms: relativistic energy, not relativistic motion}\label{sec:heavy-atoms}

The most-cited evidence that atoms are relativistic comes from the heavy elements, and it is real. Gold is
yellow where silver is white --- the relativistically narrowed $5d\to6s$ gap dropping into the visible. Mercury
is liquid where its neighbours are solid --- the contracted, inert $6s^2$ pair weakening the metallic bond. Most
of the lead--acid cell voltage vanishes if the relativistic terms are removed from the calculation.
Non-relativistic calculations get these \emph{systematically} wrong, by amounts growing as $(Z\alpha)^2$, the
signature of a relativistic kinetic correction. RealQM neither can nor should deny this.

But RealQM reads it in a way the standard picture cannot. In the standard telling the effect is ascribed to
inner electrons ``moving at a large fraction of $c$'' --- an incoherent statement for a stationary bound state,
since nothing in a stationary state moves anywhere. In RealQM there is no such motion: the ground-state electron
is a \emph{static charge density}, its velocity $\mathbf v=\nabla S/m=0$. The word ``speed'' does not apply.

What is large is not a speed but the \emph{kinetic energy of the static density itself}. Even at rest, a
localised charge density carries a kinetic energy
\begin{equation}
T=\int \frac{|\nabla\psi|^2}{2m}\,d^3x,
\label{eq:localisation-energy}
\end{equation}
the energy cost of its own curvature --- the same Heisenberg localisation energy that stabilises matter against
collapse. For a tightly Coulomb-confined inner shell at high $Z$ this energy, equivalently the momentum content
$\langle p^2\rangle=\int|\nabla\psi|^2$, is a substantial fraction of $m c^2$: $\langle p^2\rangle/m^2c^2\sim
(Z\alpha)^2$. And there the non-relativistic form \eqref{eq:localisation-energy} is simply the wrong functional;
the correct kinetic energy of the density is the relativistic one,
\begin{equation}
T_{\rm rel}=\int\Big(\sqrt{p^2c^2+m^2c^4}-mc^2\Big)\,|\tilde\psi(p)|^2\,d^3p,
\label{eq:relativistic-kinetic}
\end{equation}
a functional of the \emph{static} density's momentum distribution $|\tilde\psi(p)|^2$, with no motion implied.

So the heavy-element evidence is evidence for the \emph{form of the kinetic-energy functional} of a static
charge cloud, not for velocity. RealQM removes the incoherent orbiting-at-$0.6c$ and replaces it with a coherent
statement about the energy of a static distribution: at high $Z$ the cloud's own kinetic energy enters the regime
where the relativistic dispersion matters. The gold ring and the liquid thermometer are showing us the
kinetic-energy functional of a static charge density, not a fast particle --- a cleaner reading than the standard
one, which must speak of a motion its own stationary states forbid.

This leaves RealQM with a clear liability, and no evasion. Its kinetic term is the non-relativistic
\eqref{eq:localisation-energy}, correct for light atoms and confirmed by its ${\sim}1\%$ energies, but wrong for
heavy ones. To compute gold, mercury, or lead, RealQM must adopt the relativistic kinetic functional
\eqref{eq:relativistic-kinetic}. This is \emph{not} the importation of relativistic motion --- the electron stays
static --- but the correct energy of a static density whose localisation is extreme. RealQM does not yet carry
it; that is the real boundary, stated without the misleading language of speed.

\section{The honest boundary}\label{sec:rel-boundary}

What, then, genuinely requires relativity, and is therefore outside or at the edge of RealQM's non-relativistic
reach? Three things, named plainly (and consistent with Chapter~\ref{ch:limits}).
\begin{itemize}
\item \textbf{The magnitude of fine structure.} That spectra split is one thing; that the splitting has the
observed $v^2/c^2$ size is a relativistic fact, and RealQM does not yet deliver fine-structure magnitudes.
\item \textbf{High energy and pair creation.} Antimatter, particle creation, and the parton/deep-inelastic regime
are genuinely relativistic and outside the scope of a Coulomb continuum theory of bound charge.
\item \textbf{The precision anomaly.} The measured $g=2.00232\ldots$, the anomalous moment, is a QED effect; RealQM
reaches the geometric $g=2$, not (yet) its radiative correction.
\end{itemize}

\section{Conclusion: sufficient, not necessary}\label{sec:rel-conclusion}

The role of relativity in the physics of matter has been systematically \emph{overstated}. For the structure,
energetics, chemistry, and orbital magnetism of atoms, molecules, and nuclei, relativity does no work, and RealQM
computes these without it. The one phenomenon most often cited as forcing relativity onto the electron --- the
factor $g=2$ --- is non-relativistic geometry, the SU(2) double cover, not a consequence of Lorentz invariance.
The point electron whose kinematics would demand relativity is replaced by a charge configuration obeying
continuum mechanics. What is genuinely relativistic --- the Lorentz invariance of light, the magnitude of fine
structure, high-energy pair creation, the QED anomaly --- is either coupled in classically (Maxwell), left
honestly out of scope, or, in the case of spacetime itself, reconstructed from observation rather than assumed.
Relativity, in short, is \emph{sufficient} for much of what it is credited with and \emph{necessary} for far
less.

One point deserves emphasis, because it is easily mistaken. RealQM applies \emph{one and the same} variational
principle to every element, light and heavy alike --- there is no separate relativistic machinery bolted on at
high $Z$, as the standard theory bolts the Dirac equation onto the Schr\"odinger equation. The electrons never
move: for hydrogen and for uranium alike the ground state is a static charge density. What differs between light
and heavy elements is not the theory and not any velocity, but a single number --- the kinetic-energy content of
the confined charge, $\langle p^2\rangle/m^2c^2\sim(Z\alpha)^2$, negligible for light atoms and appreciable for
heavy ones. The one universal kinetic functional that is correct at all $Z$ is the relativistic dispersion
$\sqrt{p^2c^2+m^2c^4}-mc^2$ of \eqref{eq:relativistic-kinetic}, which reduces automatically to $p^2/2m$ where the
content is small. Adopting it changes nothing about RealQM's ontology --- static charge, no motion, no spacetime
background --- and keeps the theory uniform across the periodic table, more so than a formalism that must switch
equations at high $Z$. So the honest statement is not that the atom is relativistic or non-relativistic, but that
its electrons are always \emph{still}: what grows with $Z$ is the energy of their confinement, and RealQM's one
task there is to use the correct kinetic energy of a static charge --- an energy, never a motion.
